\documentclass[a4paper,12pt]{article}
\usepackage[spanish,activeacute]{babel}
\usepackage[latin1]{inputenc}

%% Para las columnas
\usepackage{multicol}


\usepackage{amssymb, amsmath}
\usepackage{latexsym}
\usepackage{multicol}


%% Para numeraci?n autom?tica
\newcounter{nroproblema}
\setcounter{nroproblema}{1}
\newcounter{nroapartado}
\setcounter{nroapartado}{1}
\newcommand{\n}{%
    \vspace{.1in}\noindent{\bf \arabic{nroproblema}.\ }%
    \addtocounter{nroproblema}{1}%
    \setcounter{nroapartado}{1}%
    }%
\newcommand{\apartado}{\vspace*{0.2cm}\hspace*{0.2cm}%
   {\bf \alph{nroapartado})}
   \addtocounter{nroapartado}{1}%
   }

%% m?genes

\setlength{\unitlength}{1mm} \addtolength{\voffset}{-15mm}
\addtolength{\textheight}{30mm} \addtolength{\hoffset}{-15mm}
\addtolength{\textwidth}{30mm}


\begin{document}

\begin{center}{\bf {\Large \'ALGEBRA I. HOJA 3}}\end{center}




\n Sea $\mathcal{M}_{m\times n}(\mathbb{K})$ el conjunto de  matrices $m\times n$ sobre el
cuerpo $\mathbb{K} = \mathbb{R}$ � $\mathbb{C}$. Comprobar que en $\mathcal{M}_{m\times n}(\mathbb{K})$
\begin{enumerate}
    \item La suma es asociativa.
    \item Cada matriz tiene un elemento
    opuesto (inverso aditivo).
    \item Sea $m=n$. Entonces el producto de matrices es distributivo (a derecha y a izquierda), con 
    respecto a la suma de
    matrices.
\end{enumerate}

    \n Dada la matriz
    $$
    A=\left(%
\begin{array}{cc}
  1 & 0 \\
  1 & 0 \\
\end{array}%
\right),
$$
comprobar que no existe una matriz $B$ tal que $AB=I=BA$.
  
\n Siendo $A$ y $B$ las matrices dadas a continuaci�n,
    calcular los productos $AB$ y $BA$ cuando sea posible y
    comparar los resultados.
    \begin{multicols}{2}
\begin{enumerate}
    \item [a)] $A=\left(%
\begin{array}{cc}
  1 & 1 \\
  -1 & 1 \\
\end{array}%
\right)\ B=\left(%
\begin{array}{cc}
  3 & 4 \\
  5 & 5 \\
\end{array}%
\right)$
    \item [b)] $A=\left(%
\begin{array}{cc}
  2 & 3 \\
  0 & 1 \\
\end{array}%
\right)\ B=\left(%
\begin{array}{cc}
  1 & 3 \\
  0 & 8 \\
\end{array}%
\right)$
    \item [c)] $A=\left(%
\begin{array}{ccc}
  2 & 4 & 0 \\
  5 & 0 & 3 \\
  0 & 1 & 0 \\
\end{array}%
\right)\ B=\left(%
\begin{array}{cc}
  5 & 50 \\
  1 & 60 \\
  1 & 86 \\
\end{array}
\right)$
    \item [d)] $A=\left(%
\begin{array}{ccc}
  2 & 3 & 4 \\
  4 & 3 & 2 \\
\end{array}%
\right)\ B=\left(%
\begin{array}{cc}
  1 & 2 \\
  1 & 0 \\
  4 & 4 \\
\end{array}%
\right)$
\end{enumerate}
    \end{multicols}
    
    \n �Es cierta para matrices cuadradas la relaci�n
    $(A+B)^2=A^2+2AB+B^2$? Probar la relaci�n, o suministrar un
    contraejemplo. Determinar que sucede en el caso especial de 
    las matrices $2\times 2$ de la forma 
    $\left(
\begin{array}{cc}
  u & -v \\
  v & u \\
\end{array}
\right).$
�Es el producto de tales matrices conmutativo? Si son
distintas de la matriz 0 �tienen siempre un inverso multiplicativo?

    \n Comprobar que el producto de dos matrices puede ser nulo
    sin que lo sean ninguno de los factores, hallando $AB$, siendo
    $A$ y $B$ las matrices dadas a continuaci�n:
        \begin{multicols}{2}
\begin{enumerate}
    \item [a)] $A=\left(%
\begin{array}{cc}
  1 & 1 \\
  1 & 1 \\
\end{array}
\right)\ B=\left(
\begin{array}{c}
  1 \\
  -1 \\
\end{array}
\right)$
    \item [b)] $A=\left(
\begin{array}{cc}
  3 & 6 \\
  1 & 2 \\
\end{array}
\right)\ B=\left(
\begin{array}{c}
  -2 \\
  1 \\
\end{array}
\right)$
    \item [c)] $A=\left(
\begin{array}{cc}
  5 & 6
\end{array}
\right)\ B=\left(
\begin{array}{c}
  -6 \\
  5
\end{array}
\right)$
    \item [d)] $A=\left(%
\begin{array}{cc}
  1 & 0 \\
  1 & 0  \\
\end{array}
\right)\ B=\left(
\begin{array}{cc}
  0 & 0 \\
  1 & 1 \\
\end{array}
\right).$
\end{enumerate}
    \end{multicols}
    
\n Siendo $A$, $B$ y $C$ las matrices dadas a continuaci�n,
    calcular los productos $D:=AB$, $DC$, $E:=BC$,  y $AE$.
�Es $DC =AE$?. En general, �es el producto de matirces $2\times 2$ asociativo?

    \begin{multicols}{3}
$A = \left(%
\begin{array}{cc}
  1 & 1 \\
  -1 & 1 \\
\end{array}%
\right)\ B = \left(%
\begin{array}{cc}
  3 & 4 \\
  5 & 6 \\
\end{array}%
\right)
    \ C = \left(%
\begin{array}{cc}
  2 & 3 \\
  0 & 1 \\
\end{array}%
\right).$
    \end{multicols}
    
    \n �Es el producto de matrices asociativo? No suponemos que las matrices son cuadradas,
pero si que los productos est�n bien definidos.
Por ejemplo, si $A$ es $2\times 1$, $B$ es $1\times 2$ y $C$ es $2\times 2$, �se cumple siempre
que  $(AB)C =A(BC)$?

 \n Dada la matriz
    $$
    A=\left(%
\begin{array}{cc}
  9/10 & 2/10 \\
  1/10 & 8/10 \\
\end{array}%
\right),
$$
calcular $A^2, A^4$, y  $A^8$. Una pregunta para meditar sobre ella: �existe $\lim_{n\to\infty} A^n$?.

\n Escribiendo $y = y(x)$ (de modo que $x$ es la variable
independiente), y con los datos siguientes, ha\-llar la
correspondiente recta de regresi�n $y = a x + b$.

Datos: $(x,y) = (1,1), (3,2), (4,4), (6,4), (8,5), (9,7), (11, 8),
(14, 9)$.


Soluci�n: $a = 6/11, b= 7/11$.

\

ADEMAS:

\

LADW, pp. 122-123, del 1.1 al 1.9. pp. 126-127, 2.1, 2.3-2.6, pensar sobre el 2.7.




\end{document}




\n Sea $P_2(\mathbb{R})$ el espacio de polinomios sobre $\mathbb{R}$ de grado $\le 2$, con
el  producto escalar
$(f,g):= \int_0^1 f(t) g(t) dt$. Hallar una base para el complemento ortogonal del subespacio generado por $1 + 3t$.

\n Hallar una base para el complemento ortogonal del subespacio generado por $(1,2,3)$ en $\mathbb{R}^3$. 

\n Decidir razonadamente si los siguientes conjuntos, con las  adici\'on y multiplicaci\'on por escalares definidas de manera natural, forman espacios vectoriales:
\begin{enumerate}
\item   $\{f\colon [0,1]\to \mathbb{R} \text{ tal que } f \text{ es continua}\}$
\item   $\{f\colon [0,1]\to \mathbb{R}_+ \text{ tal que } f \text{ es continua}\}$, donde $\mathbb{R}_+$ denota los n\'umeros positivos,
\item    $\{f\colon [0,1]\to \mathbb{R} \text{ tal que }  f \text{ es continua}\}$,
\item $P_n(\mathbb{R})$, el espacio de polinomios reales de grado $\le n$ (donde $n$ es un n\'umero entero positivo fijado),
\item  el espacio de polinomios reales de grado \emph{igual} a $n$.\end{enumerate}

\n Sea $P(\mathbb{R})$ el espacio de polinomios  sobre $\mathbb{R}$. Comproba que,  con las  adici\'on y multiplicaci\'on por escalares definidas de manera natural,  es una espacio vectorial. Existe una base de 
$P(\mathbb{R})$?
